\documentclass[10pt,a4paper]{amsart}
\usepackage[margin=19mm]{geometry}
\usepackage{amsmath,amssymb,mathtools}
\usepackage[scr=rsfs,cal=euler]{mathalfa}
\usepackage{xfrac}
\usepackage[unicode,hidelinks,bookmarks=false]{hyperref}
\newcommand{\ie}{i.e.\ }
\newcommand{\cf}{cf.\ }
\newcommand{\resp}{respectively\ }
\newcommand{\Subsection}{\subsection}
\providecommand{\nobreakdash}{\nobreak\hbox{-}}
\setlength{\emergencystretch}{2em}
\multlinegap0pt
\usepackage{bm}
\usepackage{bbm}
\usepackage{dsfont}
\usepackage{xspace}
\usepackage{cancel}
\usepackage{mathtools}
\usepackage[shortlabels]{enumitem}
\usepackage{amsthm}
\makeatletter
\@ifundefined{lemma}{\newtheorem{lemma}{Lemma}}{}
\@ifundefined{proposition}{\newtheorem{proposition}{Proposition}}{}
\@ifundefined{remark}{\newtheorem{remark}{Remark}}{}
\makeatother
\makeatletter
\newcommand{\R}{\mathbb R} 
\@ifundefined{pf}{\newenvironment{pf}{\it Proof. \enskip \rm}{\hfill $\qed$}}{}
\renewcommand{\phi}{\varphi}
\@ifundefined{proofof}{\newenvironment{proofof}[1]{\smallskip\noindent\emph{Proof of #1.}%
\hspace{1pt}}{\hspace{-5pt}{\nobreak\quad\nobreak\hfill\nobreak%
$\square$\vspace{8pt}\par}\smallskip\goodbreak}}{}
\makeatother
\title[Shock wavefronts for parabolic equations with sign-changing ]{Shock wavefronts for parabolic equations with sign-changing diffusivity}
\author{Diego Berti, Andrea Corli, Luisa Malaguti}
\date{Source: arXiv:2604.19595v1 (CC BY 4.0)}
\begin{document}
\maketitle

\noindent\emph{Source and licence.} Shock wavefronts for parabolic equations with sign-changing diffusivity. Version arXiv:2604.19595v1 (CC BY 4.0), distributed under the Creative Commons Attribution 4.0 International licence, as recorded in the arXiv licence metadata for this exact version and shown on the abstract page https://arxiv.org/abs/2604.19595v1. Full rights notice: licenses/arxiv-2604.19595-CC-BY-4.0.txt.\par\smallskip

\noindent\textit{Editorial scope.} Counted unit: the Lemma whose source label is \texttt{l:converse} in the article (source0.tex lines 1389--1392, source label \texttt{l:converse}), delivered together with the article's own complete original proof of it (source0.tex lines 1394--1435, one \texttt{proof} environment of 42 lines). No mathematics is written, repaired or re-derived: statement and proof are the article's own text line for line. Required context retained uncounted: e:E (lines 137--140); e:phi-limits (lines 214--216); e:EP (lines 220--222); e:solution (lines 794--796); p:Pconditions (lines 854--861); e:llsk (lines 900--902); e:hypo (lines 911--915); p:tratti (lines 928--945); e:z (lines 968--970); l:extended (lines 979--980); p:twins (lines 1050--1061); rem:2jumps (lines 1183--1205); e:gluing (lines 1342--1349). Every definition, display and label of those lines is retained unchanged. Omitted from this package and not redistributed: 1--136, 141--213, 217--219, 223--793, 797--853, 862--899, 903--910, 916--927, 946--967, 971--978, 981--1049, 1062--1182, 1206--1341, 1350--1388, 1393--1393, 1436--1715. Nothing inside a retained excerpt is omitted or altered: every retained body is delivered exactly as the article's lines are written, so no omission receipt of any kind accompanies this package. Citations: the retained excerpts carry no citation command at all, so no bibliography entry of the article is transcribed and no reference is redistributed. Reference adaptation: none was needed and none was made; every reference inside the delivered unit resolves inside it, so no unresolved reference or citation is produced. Printed slips kept literal: no printed word, symbol, hypothesis or number of the counted statement or of its proof was altered. Layout and class: the article's own class is \texttt{article} with packages including amsthm, amsfonts, amssymb, amsmath, amscd, latexsym, eucal, enumitem, graphics, graphicx, mathscinet, psfrag, tikz, footmisc; the wrapper is the lane's pinned \texttt{amsart} editorial wrapper at 10pt on A4, which restates the theorem-like environments the retained text needs and adds the article's own macro definitions that the retained text needs (\texttt{\textbackslash R}, \texttt{\textbackslash phi}), with extra wrapper packages (bm, bbm, dsfont, xspace, cancel, mathtools, enumitem). The delivered numbering is the unit's own. Assets: no retained span contains a figure environment, a graphics-inclusion command, a PDF-page inclusion, a table or a TikZ picture; no image file is copied into this package, the package declares no asset dependency and no image file is redistributed with it. Licence: version arXiv:2604.19595v1 of this article is distributed under the Creative Commons Attribution 4.0 International licence, as recorded in the arXiv licence metadata for this exact version and shown on the abstract page https://arxiv.org/abs/2604.19595v1. A full rights notice is delivered at licenses/arxiv-2604.19595-CC-BY-4.0.txt. Characters: every retained excerpt is pure ASCII, so the delivered engine, XeLaTeX with its default Latin Modern text font, reports no missing character.
\begin{equation}\label{e:E}
u_t =P(u)_{xx} + g(u), \qquad t\ge 0, \, x\in \R.
%u_t =\left(D(u)u_x\right)_x + g(u), \qquad t\ge 0, \, x\in \R.
\end{equation}

\begin{equation} \label{e:phi-limits}
\phi(-\infty)=1\quad \hbox{and}\quad \phi(\infty)=0.
\end{equation}

\begin{equation}\label{e:EP}
P(\phi)''+c\phi'+g(\phi)=0,
\end{equation}

\begin{equation}\label{e:solution}
\int_I P(\phi)\psi''\, d\xi = c\int_I \phi\psi'\, d\xi -\int_I g(\phi)\psi\, d\xi, \qquad \psi\in C_c^\infty(I).
\end{equation}

\begin{proposition}\label{p:Pconditions}
Assume {\rm (R)} and let $\phi$ be a traveling-wave solution of equation \eqref{e:E} defined on $I$ with speed $c$. Then, 
\begin{enumerate}[label=(\roman*)]
\item $P(\phi)$ is an absolutely continuous function in $I$ and $P(\phi) \in C^1(I_*)$;
\item $P(\phi)'+c\phi$ can be extended to $I_{s}$ by continuity;  
\item $\phi$ is of class $C^2$ and a classical solution of \eqref{e:EP} in ${\rm int}(I_*)$.
\end{enumerate}
\end{proposition}

\begin{equation}\label{e:llsk}
\left[P(\phi)\right]=0 \quad \text{and} \quad  \left[P(\phi)'+c\phi\right]=0 \ \mbox{ in } \ I.
\end{equation}

\begin{equation}\label{e:hypo}
P(\phi_r)-P(\phi_{\ell})=\int_{\phi_\ell}^{\phi_r}D(s)\,d s=0
\quad \mbox{ and } \quad
c=-\frac{\displaystyle P(\varphi)'(\xi_s^-)-P(\phi)'(\xi_s^+)}{\phi_r-\phi_\ell}.
\end{equation}

\begin{proposition}\label{p:tratti}  Assume {\rm (R)}. Then $\phi $ is a shock traveling-wave solution in $I$ of \eqref{e:E} with $\phi'(\xi) = 0$ a.e.\! in $I$ if and only if there exist $-\infty\le a_0<a_1<\ldots<a_m\le \infty$, with $m>1$ and $I=(a_0, a_m)$, $\gamma_k\in [0,1]$ with $\gamma_k\ne \gamma_{k+1}$ for $k=1,\ldots,m-1$, such that 
\begin{equation}\label{e:ggPP}
g(\gamma_1)=\cdots =g(\gamma_m)=0
\quad \mbox{ and } \quad 
P(\gamma_1)=\cdots=P(\gamma_m).
\end{equation}
In this case, $\phi$ is a step function of the form
\begin{equation}\label{e:step}
\phi(\xi)
= \sum_{k=1}^{m} \gamma_k \, \chi_{k}(\xi), \quad  \xi\in I, 
\end{equation}
where \begin{equation*} \chi_k(\xi)=\begin{cases} 1 \ &\mbox{ if } \ \xi \in (a_{k-1},a_k]\cap I,
\\ 
0 \ &\mbox{otherwise}
\end{cases}
\end{equation*}
and  $c=0$.
\end{proposition}

\begin{equation}\label{e:z}
z(\mu):=D(\mu)\phi'\left(\phi^{-1}(\mu)\right) \neq 0, \quad \mu \in \varphi(J).
\end{equation}

 \begin{lemma}\label{l:extended} Let $\phi$ be the profile of a monotone regular traveling-wave solution and  $z$ the corresponding function as in  \eqref{e:z}. Then $z$ can be continuously extended to $\overline{\phi(I)}$. 
 \end{lemma}

\begin{proposition}\label{p:twins}
Consider equation \eqref{e:EP}. Then, for every $c\in\R$ we have:
\begin{enumerate}[label=(\arabic*)]
\item There exists a unique regular semi-wavefront from $1$, connecting $1$ to $\beta$.

Denote by $\phi_c$ its profile and by $z_c$ the continuous extension to $[\beta,1]$ of the corresponding funtion defined in \eqref{e:z}.  If $c_1 < c_2$ then $z_{c_1}(\lambda) < z_{c_2}(\lambda)$ for every $\lambda \in (\beta, 1)$, and $z_{c_1}(\beta) \le z_{c_2}(\beta)$.

  \item There exists a unique regular semi-wavefront to $0$, connecting $\alpha$ to $0$.

Let $\psi_c$ be its profile and $z_c$ the continuous extension to $[0,\alpha]$ of the corresponding function defined in \eqref{e:z}. If $c_1 < c_2$ then $z_{c_1}(\mu) > z_{c_2}(\mu)$ for every $\mu \in (0,\alpha)$, and $z_{c_1}(\alpha)\ge z_{c_2}(\alpha)$.
\end{enumerate}
\end{proposition}

\begin{remark}
\label{rem:2jumps}
{\rm 
By Proposition \ref{p:tratti} we have that $\phi$ is a shock wavefront of \eqref{e:E}-\eqref{e:phi-limits} with  $\phi'=0$ a.e. in $\mathbb{R}$ if and only if \begin{itemize}
\item either $P(0)=P(1)$ and, for some $\xi_0 \in \mathbb{R}$,
\[
\phi(\xi)=\begin{cases}
1 & \text{for } \xi \in (-\infty,\xi_0], \\%[4pt]
0 & \text{for } \xi \in (\xi_0,\infty);
\end{cases}
\]

\item or $P(0)=P(\gamma)=P(1)$ and, for some $\xi_0^1, \xi_0^2 \in \mathbb{R}$ with $\xi_0^1 < \xi_0^2$,
\[
\phi(\xi)=\begin{cases}
1 & \text{for } \xi \in (-\infty, \xi_0^1], \\%[4pt]
\gamma & \text{for } \xi \in (\xi_0^1, \xi_0^2], \\%[4pt]
0 & \text{for } \xi \in (\xi_0^2,\infty).
\end{cases}
\]
\end{itemize}
As a consequence, $c=0$.}
\end{remark}

\begin{equation}
\label{e:gluing}
\phi(\xi) := 
\begin{cases}
\hat\phi(\xi), & \xi \in (-\infty, \xi_s], \\%[4pt]
\check\phi(\xi), & \xi \in (\xi_s, \infty),
\end{cases}
\end{equation}

\begin{lemma}
\label{l:converse}
Assume $c \in \mathbb{R}$. Let $\phi$ be defined as in \eqref{e:gluing} for some $\xi_s$ in $\mathbb{R}$. If both conditions in \eqref{e:hypo} are satisfied, then $\varphi$ is the unique profile with speed~$c$ of a shock wavefront with a single jump from $\varphi_r$ to $\varphi_\ell$ located at $\xi_s$.
\end{lemma}

\begin{proof}
Let $\phi$ be defined as in \eqref{e:gluing} for some $c, \xi_s \in \mathbb{R}$. Assume  conditions  \eqref{e:hypo}. 

If  $(\varphi_\ell,\varphi_r)=(0,1)$, then \eqref{e:hypo}$_1$ implies $P(1)=P(0)$ and \eqref{e:hypo}$_2$ yields $c=0$. By Remark \ref{rem:2jumps} it follows that $\phi$ is the unique shock traveling wave solution with a single jump. 

Now consider $(\varphi_\ell, \varphi_r)\neq (0,1)$. Since the function $\varphi$ defined in \eqref{e:gluing} is globally defined, decreasing, and continuous except at $\xi_s$, then $\phi$ is a shock wavefront if it satisfies \eqref{e:solution}. Moreover, since we already know that $\phi$ is a solution in $(-\infty,\xi_s)\cup(\xi_s,\infty)$, it suffices to show that $\phi$ also solves \eqref{e:solution} in a neighborhood $(a,b)$ of $\xi_s$, with $a,b\in \mathbb{R}$. We may further assume that $(a,b)$ is small enough not to contain the points 
$\xi_1 := \inf\{\xi:\, \varphi(\xi)<1\}$ and $\xi_2 := \sup\{ \xi:\, \varphi(\xi)>0\}$,
whenever they are finite. As a conclusion, we have $\varphi \in C^2\left((a,b)\setminus\{\xi_s\}\right)$, by Proposition~\ref{p:Pconditions}{\em (iii)}. %

Let $\psi \in C^\infty_c(\mathbb R)$, with ${\rm supp}(\psi)\subset (a,b)$, with $a$ and $b$ as above. We need to prove 
 \begin{equation}
 \label{e:intCa}
 \int_a^b P(\varphi)\psi''-c\varphi \psi' +g(\varphi)\psi\,d\xi = 0.
 \end{equation}
Denote $\mathcal L:= P(\varphi)\, \psi'' - c \varphi \,\psi' +g(\varphi)\,\psi$.
In every interval $(c,d)\subset(a,b)$ such that $\varphi \in C^2(c,d)\cap C^0[c,d]$ and $P(\varphi)' \in C^0[c,d]$  we have
\begin{align}
\int_c^d \mathcal L\,d\xi &= \int_c^d P(\varphi)\, \psi'' - c \varphi \,\psi' +g(\varphi)\,\psi\,d\xi
\nonumber
\\
&= P(\varphi)\,\psi'|^d_c + \int_c^d  -\left(P(\varphi)' +c\varphi\right)\psi' +g(\varphi)\,\psi\,d\xi 
\nonumber
\\
&= P(\varphi)\,\psi'|^d_c -\left(P(\varphi)'+c\varphi\right)\psi|^d_c+\int_c^d  \left(P(\varphi)''+c\varphi' +g(\varphi)\right)\psi\,d\xi
\nonumber
\\
&= P(\varphi)\,\psi'|^d_c -\left(P(\varphi)'+c\varphi\right)\psi|^d_c,
\label{e:heks}
\end{align}
since $P(\varphi)''+c\varphi' +g(\varphi)=0$ in $(c,d)$. By Lemma \ref{l:extended}, $P(\varphi)'$ is continuously extendable to the intervals $[a,\xi_s]$ and $[\xi_s,b]$; therefore we can apply \eqref{e:heks} to both of them. Then \eqref{e:intCa} holds if and only if
\begin{equation*}
\begin{aligned}
0&=\int_a^b \mathcal L \, d\xi  = \int_a^{\xi_s} \mathcal L \,d\xi+ \int_{\xi_s}^b \mathcal L\,d\xi
\\
&= P(\hat\varphi)\psi'|_a^{\xi_s} + P(\check\varphi)\psi'|_{\xi_s}^b
-\left(P(\hat\varphi)'+c\hat\varphi\right)\psi|_a^{\xi_s} - \left(P(\check\varphi)'+c\check\varphi\right)\psi|_{\xi_s}^b
\\
&=[P(\varphi)]\psi'(\xi_s)-\left[P(\varphi)'+c\varphi\right]\psi(\xi_s).
\end{aligned}
\end{equation*}
Since $\psi$ is arbitrary, \eqref{e:intCa} holds if and only if  \eqref{e:llsk} are satisfied. These conditions are equivalent to those in  \eqref{e:hypo}. Therefore, $\phi$ is a shock wavefront. The uniqueness (up to shifts) of $\phi$ then follows from Proposition~\ref{p:twins}.
\end{proof}

\end{document}
